\noindent\textbf{Resumo:}

\noindent O direito de acesso à informação e a transparência administrativa pressupõem que o Estado conheça e evidencie o funcionamento de seus próprios procedimentos. Contudo, sistemas eletrônicos como o Sistema Eletrônico de Informações (SEI) geram vastos rastros de atividades diárias sem consolidar o fluxo real dos processos, dificultando auditorias e a otimização da máquina pública. O problema desta pesquisa consiste em determinar se é possível extrair, a partir desses registros brutos, um catálogo auditável e confiável dos processos efetivamente executados. Trata-se de uma pesquisa aplicada e quantitativa, conduzida como estudo de caso na Secretaria de Administração do Estado do Piauí (SEAD-PI), sobre 5.243 processos concluídos em 2025. Utilizando mineração de processos e Inteligência Artificial (IA), construiu-se um grafo de conhecimento de eventos. Os processos foram agrupados combinando o texto do documento inicial ao histórico de tramitação. Para ajustar o caminho típico de cada processo, foram avaliadas seis famílias de modelos, incluindo cadeias de Markov e redes neurais recorrentes (como a arquitetura GRU), testadas em três alfabetos de traço (ação por unidade, roteamento e procedimento). Os resultados revelam um ganho positivo e consistente em 252 comparações contra partições-controle embaralhadas. Na modelagem global, observou-se eficácia preditiva no alfabeto de procedimento (ganho de 0,109 com Markov de terceira ordem e 0,131 com redes neurais GRU). Embora as métricas internas de clusterização não tenham isolado um único agrupamento ideal, a aplicação destas técnicas gerou uma entrega concreta: um catálogo validado de 60 tipos de processos com cobertura de 75,1% no alfabeto de roteamento, preservando as margens de anonimização exigidas. Conclui-se que o uso de IA sobre logs administrativos supre a lacuna empírica dos registros, viabilizando a governança orientada a dados.

\vspace{0.6\baselineskip}
\noindent\textbf{Palavras-chave:} \textit{mineração de processos; inteligência artificial; redes neurais; otimização; transparência administrativa.}.

\vspace{1.2\baselineskip}
\noindent\textbf{Abstract:}

\noindent Access-to-information law presumes a record of how administrative procedures run, yet
empirical work on open-records requests finds that the binding obstacle is not refusal but the fact
that this record was never created. The problem is whether an auditable catalogue of the process
types an administration actually executes, and of the typical path of each, can be derived from the
administration's own event log, and on what evidence it may be considered trustworthy. This is
applied, quantitative research conducted as a single case study at the Secretariat of
Administration of the State of Piaui, over the 5,243 cases concluded during 2025 in the Electronic
Information System: an event knowledge graph is built, cases are grouped by opening-document text
together with routing and action history, and the typical path of each type is fitted by six model
families within each of three trace alphabets --- action-by-unit, routing and procedure ---,
against three controls. The gain over the shuffled control partition, which preserves cluster
count, size distribution and evaluated population, is positive in all 252 comparisons of the
lattice, in all three alphabets. Over a single global model, with the model family held fixed,
there is an effect only in the procedure alphabet --- 0.109 for a third-order Markov chain and
0.131 for the GRU ---, whereas in action-by-unit and in routing the measured gain is
indistinguishable from the residual of the evidence load itself. Accuracy separates the fourteen
candidate clusterings by 0.0055, and none of the eight internal indices tested predicts downstream
quality in a sample that resolves only $|\rho| \geq 0.53$: there is no best clustering to elect.
The delivered catalogue has 60 types and coverage of 0.751 in the routing alphabet; the
anonymisation the authorising instrument requires is non-inferior at the pre-registered margin in
all 54 comparisons.

\vspace{0.6\baselineskip}
\noindent\textbf{Keywords:} \textit{process mining; clustering; graphs;
administrative transparency}.

\newpage
